\documentclass[11pt]{article}

\usepackage[margin=1in]{geometry}
\usepackage[T1]{fontenc}
\usepackage[utf8]{inputenc}
\usepackage{amsmath,amssymb,array}
\usepackage[hidelinks]{hyperref}

\newcommand{\IUT}{IUT}
\newcommand{\Lean}{Lean}
\newcommand{\R}{\mathbb{R}}
\newcommand{\F}{\mathbb{F}}
\newcommand{\ThetaPilot}{\Theta\text{-pilot}}
\newcommand{\qPilot}{q\text{-pilot}}

\title{A First Lean Formalization Audit of the IUT III Corollary 3.12 Corridor}
\author{IUT Lean formalization project}
\date{May 28, 2026}

\begin{document}
\maketitle
\sloppy

\begin{abstract}
This note records the present state of a Lean formalization of a narrow part of
Mochizuki's inter-universal Teichmuller theory: the passage from IUT III,
Theorem 3.11 to Corollary 3.12, with attention to the Scholze--Stix objection.
The current Lean code proves conditional route theorems and obstruction lemmas.
It does not prove Corollary 3.12 and does not refute it.  The main formal result
so far is that a nonarchimedean \((\mathrm{Ind3})\) entry, together with explicit
Hodge/SHE square-weight transport data, feeds a hull/log-volume
\(q\)-to-\(\Theta\) comparison, while several simplified collapse mechanisms are
rejected by Lean.
\end{abstract}

\section{Scope}

The review used the local Markdown conversions of IUT I--IV, Mochizuki's 2026
formalization note, and the Scholze--Stix note.  The relevant source passages are:
\begin{itemize}
\item IUT I: initial \(\Theta\)-data, prime strips, Hodge theaters, and
  non-scheme-theoretic \(\Theta\)-links.
\item IUT II: Hodge--Arakelov evaluation, theta values of shape \(q^{j^2}\),
  multiradiality, and Gaussian monoids.
\item IUT III: Theorem 3.11; Corollary 3.12, especially Steps (x) and (xi);
  Remarks 3.12.2 and 3.12.3.
\item IUT IV: later explicit log-volume estimates and Diophantine consequences.
\item Scholze--Stix: the ordered real-line identification problem in their
  Section 2.2.
\item Mochizuki 2026: the ``3.11.5 \(\Rightarrow\) 3.12'' fourth-triangle
  decomposition.
\end{itemize}

\section{Mathematical Target}

IUT III Step (x) treats procession-normalized log-volumes as invariants for
\((\mathrm{Ind1})\) and \((\mathrm{Ind2})\), and turns \((\mathrm{Ind3})\) into
an upper inequality.  Step (xi) then uses the multiradial construction, IPL, SHE,
holomorphic hulls, determinants, and log-volume to place the \(q\)-pilot
log-volume in an upper ray controlled by the \(\Theta\)-pilot log-volume:
\[
  -|\log(q)| \in \R_{\le -|\log(\Theta)|}
  \quad\Rightarrow\quad
  -|\log(q)| \le -|\log(\Theta)|.
\]
Mochizuki's 2026 note isolates the final comparison as:
\[
\begin{array}{ccccc}
\text{Theorem 3.11}
& \xrightarrow{\mathrm{APT}+\mathrm{IPL}}
& \text{``3.11.5''}
& \xrightarrow{\mathrm{SHE}+\mathrm{IPL}}
& \text{Corollary 3.12}.
\end{array}
\]
The present formalization targets only this corridor, not the full construction
of Theorem 3.11.

\section{Dispute Interface}

Scholze--Stix require explicit bookkeeping for all ordered one-dimensional real
vector spaces.  Their pressure point is that the \(\Theta\)-side concrete
objects carry \(j^2\)-scaled degrees, while a consistent identification of the
real lines appears to leave no room for inserting all these scalars.

The formalization therefore separates:
\[
\begin{array}{c}
\text{representative } j^2 \text{ data}\\
\text{balanced sign-compatible data}\\
\text{aggregate averaged data}\\
\text{final hull/log-volume data}
\end{array}
\qquad
\text{and}
\qquad
\begin{array}{c}
\text{raw real equality}\\
\text{transported equality}\\
\text{matched transport scales}\\
\text{cancellable ordered-line comparison.}
\end{array}
\]
This separation is the main correctness guard in the current code.

\section{Lean Coverage}

The relevant implementation is in
\[
\texttt{Iut/Stage1/IUTStage1Source.lean}
\quad\text{and}\quad
\texttt{Iut/Stage1/IUTStage1Experiments.lean}.
\]
The code currently covers:
\begin{itemize}
\item labeled real-line copies and positive-scale transports;
\item a finite \(\F_\ell\)-label model with representative square weights
  \(j \mapsto j^2\);
\item a representative pilot-degree profile
  \(\deg(\Theta_j)=j^2\deg(q)\);
\item proof that one label-independent scale cannot absorb all representative
  \(j^2\) scales in this model;
\item proof that, for nonzero \(\deg(q)\), one label-independent scale cannot
  account for all representative \(\Theta_j\)-degrees;
\item proof that the raw representative \(j^2\) theta-degree rule does not
  descend to the sign quotient \(|F_\ell|=F_\ell/\{\pm1\}\) when \(\deg(q)\ne0\);
\item an absolute-label exponent on \(|F_\ell|=\{0\}\cup F_\ell^\times/\{\pm1\}\)
  which agrees with \(j^2\) for the canonical representatives
  \(0\le j\le(\ell-1)/2\), and the corresponding theta-degree profile;
\item the degree-level Gaussian evaluation
  \(q\mapsto(q^{j^2})_{j\in |F_\ell|}\), including identity at \(j=1\);
\item comparison levels:
  pointwise representative, aggregate representative, balanced sign-compatible,
  and hull/log-volume;
\item endpoint audits for the \(3.11.5 \Rightarrow 3.12\) comparison chart;
\item nonarchimedean and archimedean \((\mathrm{Ind3})\) local orientations;
\item ordered real-line alignment with cancellation only after scale equality;
\item a canonical unit-scale ordered alignment from a nonarchimedean
  \((\mathrm{Ind3})\) entry;
\item a bridge from factored square/full-label SHE obligations to the
  Hodge-descent \((\mathrm{Ind3})\) route.
\end{itemize}

\section{Current Lean Results}

The current first-pass experiment report evaluates to:
\[
\begin{array}{ll}
\texttt{orderedRealLineRouteTheoremAvailable} & = \texttt{true},\\
\texttt{nonarchimedeanEntryCanonicalAlignmentTheoremAvailable} & = \texttt{true},\\
\texttt{factoredSHEBridgeTheoremAvailable} & = \texttt{true},\\
\texttt{mismatchCounterexampleBlocksRawCancellation} & = \texttt{true},\\
\texttt{labelIndependentJ2CollapseRejectedInZModModel} & = \texttt{true},\\
\texttt{representativeJ2SignQuotientDescentRejectedInZModModel} & = \texttt{true},\\
\texttt{absLabelThetaDegreeHalfRangeModelAvailable} & = \texttt{true},\\
\texttt{gaussianDegreeEvaluationTheoremAvailable} & = \texttt{true},\\
\texttt{balancedLevelRejectedAtFinalRouteTheoremAvailable} & = \texttt{true},\\
\texttt{disputeSettledByCurrentStage} & = \texttt{false}.
\end{array}
\]
The theorem-level route is:
\[
\begin{array}{ccccc}
\text{nonarch. }(\mathrm{Ind3})\text{ entry}
& \to &
\text{canonical ordered } \R\text{-alignment}
& \to &
\text{source/target alignment}
\\[3pt]
& & \downarrow & & \downarrow
\\[-2pt]
\text{factored SHE transport}
& \to &
\text{Hodge-descent audit}
& \to &
q_{\mathrm{signed}}\le \Theta_{\mathrm{signed}}.
\end{array}
\]
The route remains conditional on the displayed hypotheses.

\section{Audit Findings}

No formal statement currently proves IUT III, Corollary 3.12 from Mochizuki's
published hypotheses.  The strongest positive result is conditional: if the
nonarchimedean \((\mathrm{Ind3})\) entry, ordered real-line alignment, and
Hodge/SHE square-weight transport audit are supplied, then Lean checks the final
signed \(q\)-to-\(\Theta\) inequality.

The strongest negative result is also limited: in the current \(\F_\ell\)-label
model, Lean rejects a single label-independent scale accounting for all
representative \(j^2\) factors, and it rejects balanced sign-compatible evidence
as the final hull/log-volume comparison level.  Lean also rejects direct descent
of arbitrary raw representative \(j^2\) theta-degree data to \(|F_\ell|\), while
accepting the half-range exponent convention stated in IUT II.

One correctness issue was fixed during this audit.  The experiment report fields
were renamed so that they describe theorem availability under hypotheses rather
than an unconditional mathematical fact.

\section{Missing Mathematics}

The present code does not yet formalize:
\begin{itemize}
\item initial \(\Theta\)-data as in IUT I;
\item actual Hodge theaters, Frobenioids, prime strips, log-shells, or
  \(\Theta\)-links;
\item the Hodge--Arakelov theta evaluation and Gaussian monoids of IUT II;
\item the construction of the actual IUT II theta values, Gaussian monoids, and
  Frobenioids behind the present degree-level Gaussian evaluation;
\item the full multiradial construction of IUT III, Theorem 3.11;
\item the genuine holomorphic hull and determinant operations of Step (xi);
\item the exact IUT IV log-volume estimates and Diophantine applications;
\item a proof that the current Hodge/SHE transport obligations follow from
  Mochizuki's full hypotheses.
\end{itemize}

\section{Next Work}

The next significant formal step is not another wrapper.  It is to replace
current explicit Hodge/SHE transport assumptions by constructions closer to the
papers:
\[
\begin{array}{c}
\text{theta values }q^{j^2}\\
\leadsto\ \text{Gaussian monoids}\\
\leadsto\ \text{square/full-label preservation}\\
\leadsto\ \text{hull/log-volume route}.
\end{array}
\]
Only after this bridge is formalized can the first-pass experiment begin to test
the actual mathematical load-bearing part of Mochizuki's response to
Scholze--Stix.

\begin{thebibliography}{9}
\bibitem{IUT1}
S. Mochizuki, \emph{Inter-universal Teichmuller Theory I:
Construction of Hodge Theaters}.

\bibitem{IUT2}
S. Mochizuki, \emph{Inter-universal Teichmuller Theory II:
Hodge-Arakelov-theoretic Evaluation}.

\bibitem{IUT3}
S. Mochizuki, \emph{Inter-universal Teichmuller Theory III:
Canonical Splittings of the Log-theta-lattice}.

\bibitem{IUT4}
S. Mochizuki, \emph{Inter-universal Teichmuller Theory IV:
Log-volume Computations and Set-theoretic Foundations}.

\bibitem{MochizukiLean}
S. Mochizuki, \emph{On the Formalization of IUT: Preliminary Progress Report},
2026.

\bibitem{ScholzeStix}
P. Scholze and J. Stix, \emph{Why abc is still a conjecture}.
\end{thebibliography}

\end{document}
